\chapter{Asset Prices: A First Model}\label{ch:b2:asset-prices}

A unit of wealth put into American shares at the end of 1949, with every dividend reinvested, was
% figdata: summary (usa_eq_final 177.6, usa_bill_final 2.37)
worth 177.6 units of the same purchasing power in 2020. Put into short-term government bills, it was
worth 2.37. The four panels below tell the same story in four countries: shares far ahead, government
bonds second, bills last. Investors who saw these returns coming would have bought shares until the gap
closed. This chapter asks why it did not close: what an asset's price should be, what return it should
pay for its risk, and how much of the gap a model of risk-averse savers can explain.

\begin{openingchart}{Cumulated real total returns, 1950--2020, in four countries: the value of one unit
  of real wealth invested at the end of 1949, all income reinvested, for equities, long-term government
  bonds and short-term bills. Times-ten vertical scale.}{Jord\`a--Schularick--Taylor Macrohistory
  Database: equity and bond total returns, bill rates, consumer prices}
\begin{tikzpicture}
\begin{groupplot}[group style={group size=2 by 2, horizontal sep=1.2cm, vertical sep=1.3cm},
    omrecord, ymode=log, log basis y=10, log ticks with fixed point, ymin=0.3, ymax=300,
    ytick={1,10,100}, xmin=1949, xmax=2020, xtick={1950,1985,2020},
    x tick label style={/pgf/number format/1000 sep=}, width=0.5\linewidth, height=0.4\linewidth,
    title style={font=\small, yshift=-1ex}]
  \nextgroupplot[title={United States}, legend to name=b2assetlegend, legend columns=3]
    \addplot[omDef, very thick] table[x=year, y=USA_equity]
      {figdata/bachelor-2/04-asset-prices/out/cumulative.csv};
    \addplot[omThm, thick] table[x=year, y=USA_bond]
      {figdata/bachelor-2/04-asset-prices/out/cumulative.csv};
    \addplot[black!60, thick, densely dashed] table[x=year, y=USA_bill]
      {figdata/bachelor-2/04-asset-prices/out/cumulative.csv};
    \legend{Equities, Bonds, Bills}
  \nextgroupplot[title={United Kingdom}]
    \addplot[omDef, very thick] table[x=year, y=GBR_equity]
      {figdata/bachelor-2/04-asset-prices/out/cumulative.csv};
    \addplot[omThm, thick] table[x=year, y=GBR_bond]
      {figdata/bachelor-2/04-asset-prices/out/cumulative.csv};
    \addplot[black!60, thick, densely dashed] table[x=year, y=GBR_bill]
      {figdata/bachelor-2/04-asset-prices/out/cumulative.csv};
  \nextgroupplot[title={France}]
    \addplot[omDef, very thick] table[x=year, y=FRA_equity]
      {figdata/bachelor-2/04-asset-prices/out/cumulative.csv};
    \addplot[omThm, thick] table[x=year, y=FRA_bond]
      {figdata/bachelor-2/04-asset-prices/out/cumulative.csv};
    \addplot[black!60, thick, densely dashed] table[x=year, y=FRA_bill]
      {figdata/bachelor-2/04-asset-prices/out/cumulative.csv};
  \nextgroupplot[title={Japan}]
    \addplot[omDef, very thick] table[x=year, y=JPN_equity]
      {figdata/bachelor-2/04-asset-prices/out/cumulative.csv};
    \addplot[omThm, thick] table[x=year, y=JPN_bond]
      {figdata/bachelor-2/04-asset-prices/out/cumulative.csv};
    \addplot[black!60, thick, densely dashed] table[x=year, y=JPN_bill]
      {figdata/bachelor-2/04-asset-prices/out/cumulative.csv};
\end{groupplot}
\node[anchor=north] at ($(group c1r2.south)!0.5!(group c2r2.south) + (0,-0.9cm)$)
  {\pgfplotslegendfromname{b2assetlegend}};
\end{tikzpicture}
\end{openingchart}

The gap is no accident of four countries. Across the 15 countries of the database with complete
% figdata: summary (n_countries 15, median_eq_mean 8.74, median_bill_mean 0.97, median_premium 7.75)
series, the median country's equities returned 8.74 per cent a year after inflation on average, its
bills 0.97 per cent: a median gap of 7.75 percentage points. But equities also lost a quarter or more of
their value in some years, which bills never did. The theory of this chapter prices that difference in
risk, and finds it far too small to account for the gap.

\begin{recall}
From \cref{ch:b2:intertemporal-choice}: the \omterm{def:b2:intertemporal-choice:discount-factor}{discount factor} $\beta$ and the \omterm{def:b2:intertemporal-choice:euler-equation}{Euler equation}, $u'(c_0) =
\beta(1 + r)u'(c_1)$. From \cref{ch:b2:insurance-and-risk-sharing}: states of the world, contingent
goods and Arrow securities; the \omterm{def:b2:insurance-and-risk-sharing:portfolio-share}{portfolio share} of a risky asset. From the Year 1 volume: arbitrage,
buying cheap in one place and selling dear in another. From the school volumes: present value.
\end{recall}

\begin{notation}[Symbols of this chapter]\label{not:b2:asset-prices:symbols}
States $s$ have probabilities $\mu_s$. Security $j$ costs $p_j$ today and pays $x_j(s)$ in state $s$; a
portfolio holds $h_j$ units of security $j$. $\psi_s$ is the \omterm{def:b2:asset-prices:state-price}{state price} of state $s$ and
$\tilde\zeta$ the \omterm{def:b2:asset-prices:sdf}{stochastic discount factor}. Gross returns are $\tilde R = \tilde x/p$; $R_f$ is the
gross risk-free return, $\tilde\xi = \tilde R - R_f$ an \omterm{def:b2:asset-prices:excess-return}{excess return}. $\beta$ is the \omterm{def:b2:intertemporal-choice:discount-factor}{discount factor};
with an asset subscript, $\beta_i$ is the asset's \omterm{def:b2:asset-prices:beta}{beta} (\cref{def:b2:asset-prices:beta}), a regression
coefficient. $\theta$ is a \omterm{def:b2:insurance-and-risk-sharing:portfolio-share}{portfolio share}, $\gamma$ the CRRA coefficient, $g$ gross consumption growth.
\end{notation}

\section{No arbitrage}\label{sec:b2:asset-prices:no-arbitrage}

Take two dates and two states of the world tomorrow, $s = 1, 2$. A security is a list of payoffs, one
per state, with a price today. A portfolio $(h_1, \dots, h_J)$ costs $\sum_j h_j p_j$ and pays $\sum_j
h_j x_j(s)$ in state $s$; a negative $h_j$ is a short sale.

\begin{definition}[No-arbitrage pricing]\label{def:b2:asset-prices:no-arbitrage}
An \emph{arbitrage} is a portfolio that costs nothing or less today and pays nothing or more in every
state, with a strict gain somewhere: a negative price, or a positive payoff in some state. Prices satisfy
\emph{no-arbitrage pricing}\index{no-arbitrage pricing} when no portfolio is an arbitrage.
\end{definition}

\begin{definition}[Law of one price]\label{def:b2:asset-prices:law-of-one-price}
Prices obey the \emph{law of one price}\index{law of one price} if any two portfolios with the same
payoff in every state have the same price.
\end{definition}

A violation of the \omterm{def:b2:asset-prices:law-of-one-price}{law of one price} is an arbitrage: buy the cheaper portfolio, sell the dearer one,
pocket the difference today and owe nothing tomorrow. So no arbitrage implies the \omterm{def:b2:asset-prices:law-of-one-price}{law of one price}. The
converse fails: a security that pays one in every state and costs $-1$ obeys the \omterm{def:b2:asset-prices:law-of-one-price}{law of one price} on
its own and is still an arbitrage.

\begin{definition}[State price]\label{def:b2:asset-prices:state-price}
The \emph{state price}\index{state price} $\psi_s$ of state $s$ is the price today of the \omterm{def:b2:insurance-and-risk-sharing:arrow-security}{Arrow security}
for state $s$: one unit of money in state $s$, nothing in the others.
\end{definition}

\begin{proposition}[State prices]\label{prop:b2:asset-prices:state-prices}
Let two securities have payoffs that are not proportional. Under the \omterm{def:b2:asset-prices:law-of-one-price}{law of one price} there are unique
numbers $\psi_1, \psi_2$ such that every portfolio's price is
\[
p = \psi_1 x(1) + \psi_2 x(2).
\]
Prices admit no arbitrage if and only if $\psi_1 > 0$ and $\psi_2 > 0$.
\end{proposition}

\begin{proof}
The two payoff vectors are linearly independent, so every payoff, in particular each \omterm{def:b2:insurance-and-risk-sharing:arrow-security}{Arrow security},
is a portfolio of the two; by the \omterm{def:b2:asset-prices:law-of-one-price}{law of one price} each has a single price, $\psi_1$ and $\psi_2$.
Every payoff $x$ is the portfolio of $x(1)$ units of the first \omterm{def:b2:insurance-and-risk-sharing:arrow-security}{Arrow security} and $x(2)$ of the second,
so its price is $\psi_1x(1) + \psi_2x(2)$. If both \omterm{def:b2:asset-prices:state-price}{state prices} are positive, a payoff that is never
negative and positive somewhere has a positive price, and a payoff of zero costs zero: there is no
arbitrage. The converse is \cref{exo:b2:asset-prices:10}.
\end{proof}

\begin{omfigure}
\begin{tikzpicture}[econ]
\begin{axis}[omrecord, grid=none, xmin=0, xmax=180, ymin=0, ymax=180, xtick={80,105,130},
    ytick={80,105,130}, xlabel={Payoff in a boom, $x(1)$}, ylabel={Payoff in a slump, $x(2)$},
    width=0.6\linewidth, height=0.6\linewidth]
  \addplot[guide, domain=0:180, samples=2] {x};
  \addplot[omDef, thick, domain=0:180, samples=2] {(100 - 0.37619*x)/0.57619};
  \addplot[omThm, thick, domain=0:180, samples=2] {(95 - 0.37619*x)/0.57619};
  \node[eqpoint] at (axis cs:105,105) {};
  \node[anchor=south west, font=\footnotesize] at (axis cs:105,105) {bond};
  \node[eqpoint] at (axis cs:130,80) {};
  \node[anchor=south west, font=\footnotesize] at (axis cs:130,80) {equity};
  \node[anchor=south west, font=\footnotesize, omDef] at (axis cs:62,136) {price 100};
  \node[anchor=north east, font=\footnotesize, omThm] at (axis cs:56,126) {price 95};
  \node[anchor=south east, font=\footnotesize, black!60, rotate=45] at (axis cs:170,170) {riskless};
\end{axis}
\end{tikzpicture}

{\small \omterm{def:b2:asset-prices:state-price}{State prices} in the payoff plane. A bond pays 105 in both states and costs 100; an equity pays
130 in a boom and 80 in a slump and costs 95 (the weekend problem's two securities). The lines join the
payoffs that cost the same; their slope is $-\psi_1/\psi_2 = -0.653$. No arbitrage means that every line
of equal price slopes down: a payoff larger in both states must cost more.}
\end{omfigure}

\section{A price as an expected discounted payoff}\label{sec:b2:asset-prices:sdf}

Divide each \omterm{def:b2:asset-prices:state-price}{state price} by the probability of its state: $\zeta_s = \psi_s/\mu_s$. Then
\[
p = \sum_s \psi_s x(s) = \sum_s \mu_s\,\zeta_s\,x(s) = \mathbb{E}[\tilde\zeta\,\tilde x].
\]

\begin{definition}[Stochastic discount factor]\label{def:b2:asset-prices:sdf}
A \emph{stochastic discount factor}\index{stochastic discount factor} is a random variable
$\tilde\zeta$, one value per state, such that every asset's price is $p = \mathbb{E}[\tilde\zeta\,\tilde
x]$; equivalently, every gross return satisfies $\mathbb{E}[\tilde\zeta\,\tilde R] = 1$. It is positive
in every state exactly when there is no arbitrage.
\end{definition}

The \omterm{def:b2:intertemporal-choice:discount-factor}{discount factor} is ``stochastic'' because a unit of money is worth more in some states than in
others. The saver of \cref{ch:b2:intertemporal-choice} says which.

\begin{proposition}[The consumption discount factor]\label{prop:b2:asset-prices:consumption-sdf}
A saver with expected utility $u(c_0) + \beta\,\mathbb{E}[u(\tilde c_1)]$ who can buy any amount of an
asset at price $p$ holds it, at an interior optimum, only if
\[
p = \mathbb{E}\Bigl[\beta\,\frac{u'(\tilde c_1)}{u'(c_0)}\,\tilde x\Bigr]:
\]
$\tilde\zeta = \beta u'(\tilde c_1)/u'(c_0)$ is a \omterm{def:b2:asset-prices:sdf}{stochastic discount factor}. The gross risk-free return
is $R_f = 1/\mathbb{E}[\tilde\zeta]$; with CRRA utility and certain growth $g$, $R_f = g^{\gamma}/\beta$.
\end{proposition}

\begin{proof}
Buying $\varepsilon$ more units costs $\varepsilon p\,u'(c_0)$ in utility today and adds
$\varepsilon\beta\,\mathbb{E}[u'(\tilde c_1)\tilde x]$ tomorrow, to first order; at an optimum the
difference is zero for small $\varepsilon$ of either sign. A risk-free asset pays $1$ in every state, so
its price is $\mathbb{E}[\tilde\zeta]$ and its gross return $1/\mathbb{E}[\tilde\zeta]$. With $u'(c) =
c^{-\gamma}$ and $\tilde c_1 = gc_0$, $\tilde\zeta = \beta g^{-\gamma}$.
\end{proof}

This is the \omterm{def:b2:intertemporal-choice:euler-equation}{Euler equation} again, one asset at a time. A unit of money is worth most in the states where
consumption is low, so $\tilde\zeta$ is high in bad times. With certain growth, the formula for $R_f$
already contains a first puzzle: with $\beta$ near one and growth of two per cent a year, a $\gamma$ of 5
asks for a real safe rate near 10 per cent, while bills have paid about one per cent
(\cref{exo:b2:asset-prices:2}).

The same formula prices bonds of every maturity. A bond paying one unit in $n$ years costs $p^{(n)} =
\mathbb{E}[\tilde\zeta_1\tilde\zeta_2\cdots\tilde\zeta_n]$, and its yield is $(p^{(n)})^{-1/n} - 1$.
If the \omterm{def:b2:intertemporal-choice:discount-factor}{discount factors} of successive years were independent, the $n$-year yield would be close to the
average of the one-year rates expected over the $n$ years: this is the \emph{expectations hypothesis}
of the term structure. It ignores any premium that a long bond pays for the risk that its price falls
when rates rise; tests of the hypothesis ask whether such a premium exists and whether it changes over
time. \Cref{ch:b2:three-equation-model} returns to yields as the channel of monetary policy.

\section{A risk premium from covariance}\label{sec:b2:asset-prices:premium}

\begin{definition}[Excess return]\label{def:b2:asset-prices:excess-return}
The \emph{excess return}\index{excess return} of an asset is its return minus the risk-free return,
$\tilde\xi_i = \tilde R_i - R_f$; its mean $\mathbb{E}[\tilde\xi_i]$ is the asset's \omterm{def:b2:expected-utility:risk-premium}{risk premium}.
\end{definition}

\begin{proposition}[Risk premium]\label{prop:b2:asset-prices:covariance}
For every asset,
\[
\mathbb{E}[\tilde R_i] - R_f = -R_f\,\mathrm{cov}(\tilde\zeta, \tilde R_i).
\]
An asset that pays well when the \omterm{def:b2:intertemporal-choice:discount-factor}{discount factor} is high, in bad times, has a return below the
risk-free return; one that pays well in good times must offer more.
\end{proposition}

\begin{proof}
$1 = \mathbb{E}[\tilde\zeta\tilde R_i] = \mathbb{E}[\tilde\zeta]\,\mathbb{E}[\tilde R_i] +
\mathrm{cov}(\tilde\zeta, \tilde R_i)$, and $\mathbb{E}[\tilde\zeta] = 1/R_f$. Multiply by $R_f$ and
rearrange.
\end{proof}

Only covariance with the \omterm{def:b2:intertemporal-choice:discount-factor}{discount factor} is priced. An asset whose return is risky but unrelated to
$\tilde\zeta$ earns the risk-free return, however volatile it is: its risk can be diversified away and
asks for no reward. With CRRA utility and small changes, $\tilde\zeta = \beta\tilde g^{-\gamma} \approx
\beta(1 - \gamma(\tilde g - 1))$, so
\[
\mathbb{E}[\tilde R_i] - R_f \approx \gamma\,\mathrm{cov}(\tilde g, \tilde R_i):
\]
the premium is the \omterm{def:b2:expected-utility:rra}{coefficient of relative risk aversion} times the covariance of the asset's return with
consumption growth. \Cref{exo:b2:asset-prices:11} derives the exact version with lognormal consumption.

\begin{definition}[Equity premium]\label{def:b2:asset-prices:equity-premium}
The \emph{equity premium}\index{equity premium} is the mean \omterm{def:b2:asset-prices:excess-return}{excess return} of a broad index of shares
over a risk-free asset, usually short-term government bills.
\end{definition}

\begin{record}[The equity premium puzzle]
Mehra and Prescott (1985) put American data for 1889--1978 into the consumption model. Over those
% ledger: b2:mehra-prescott1985
years the mean gross return on equity was $1.0698$, the gross risk-free return $1.008$, and the \omterm{def:b2:asset-prices:equity-premium}{equity
premium} $0.0618$, while consumption grew by a factor of $1.018$ a year on average with a standard
deviation of only $0.036$. Consumption that smooth makes the covariance in the premium formula small.
Taking the evidence to put the \omterm{def:b2:expected-utility:rra}{coefficient of relative risk aversion} below 10, they set it at that
% ledger: b2:mehra2003-calib
bound with $\beta = 0.99$: the model then gives a risk-free rate of 12.7 per cent and a premium of 1.4
points, against 0.8 per cent and 6.18 points in the data (as restated in Mehra 2003). Updated to 2000,
% ledger: b2:mehra2003-us
the American premium over the relatively riskless asset was 6.9 percentage points.
\end{record}

The listing below repeats the calculation on the JST data for the United States, 1950--2020. It
computes the premium, the covariance of equity returns with consumption growth, and the \omterm{def:b2:expected-utility:rra}{coefficient of
relative risk aversion} that would make the two match. It also computes the smallest $\gamma$ that could
match them if returns and consumption were perfectly correlated: since a correlation cannot exceed one,
$\mathrm{cov}(\tilde g, \tilde R) \le \sigma_g\sigma_R$, and $\gamma \ge (\mathbb{E}[\tilde\xi]/\sigma_R)/
\sigma_g$.

\omcode{code/bachelor-2/04-asset-prices/listings/puzzle.py}{1}{30}{The \omterm{def:b2:asset-prices:equity-premium}{equity premium} and the \omterm{def:b2:expected-utility:risk-aversion}{risk
aversion} it needs, United States, 1950--2020.}\label{lst:b2:asset-prices:puzzle}

\omoutput{code/bachelor-2/04-asset-prices/outputs/puzzle.txt}{A premium of 7.6 points against a
consumption growth with a standard deviation under 2 per cent.}\label{out:b2:asset-prices:puzzle}

A premium of 7.6 points needs a \omterm{def:b2:expected-utility:rra}{coefficient of relative risk aversion} of 312 at the measured
% code: listings/puzzle.py
correlation of 0.08, and at least 25 even if the correlation were one. A household with $\gamma = 25$
would pay almost a fifth of its wealth to avoid a fifty-fifty gamble of losing or gaining a fifth
(\cref{exo:b2:asset-prices:12}). The measured premium is large, the risk in aggregate consumption is
small, and the model ties the two through a single $\gamma$ that must also fit the risk-free rate.

\begin{model}{Consumption-based asset pricing}
\assumes A representative saver with time-separable CRRA utility prices every asset by its \omterm{def:b2:intertemporal-choice:euler-equation}{Euler
equation}; returns and consumption are measured without error.
\predicts Premiums proportional to the covariance of returns with consumption growth, with a factor
$\gamma$ that also fixes the risk-free rate and the saver's attitude to gambles.
\tested Mehra and Prescott (1985) on American data; the JST data for 1950--2020 (this chapter's
listing).
\breaks The \omterm{def:b2:asset-prices:equity-premium}{equity premium} needs a $\gamma$ of 25 or more (the \omterm{def:b2:asset-prices:equity-premium}{equity premium} puzzle), and such a
$\gamma$ predicts a risk-free rate far above the one observed. Models that change the utility
function, add rare disasters or drop the representative saver are the subject of the quant series and
of graduate courses.
\end{model}

\begin{measure}[How a total-return index is built]
A total-return index assumes that every dividend or coupon is reinvested in the same asset on receipt.
Its annual nominal return is the price change plus the income, $(p_t + d_t)/p_{t-1} - 1$; the real
return divides by inflation, $(1 + \text{nominal})/(1 + \pi_t) - 1$, with $\pi_t$ the change in the
consumer price index; and the index chains the returns, $V_t = V_{t-1}(1 + r_t)$. Three cautions. The
index ignores taxes and trading costs, which no investor escapes. Its composition follows the shares
the compilers could find, which tilts long histories towards companies that survived. And annual data
hide the path within the year: the worst year in the index is milder than the worst twelve months.
\end{measure}

\section{Mean--variance}\label{sec:b2:asset-prices:mean-variance}

A second route to \omterm{def:b2:expected-utility:risk-premium}{risk premiums} starts from the investor's portfolio instead of consumption. Suppose
investors care only about the mean and the variance of their portfolio's return. With two risky assets
in shares $\theta$ and $1 - \theta$, the mean is linear in $\theta$ and the variance quadratic:
\[
\mathbb{E}[\tilde R_\theta] = \theta\,\mathbb{E}[\tilde R_1] + (1 - \theta)\,\mathbb{E}[\tilde R_2], \quad
\sigma^2_\theta = \theta^2\sigma_1^2 + (1 - \theta)^2\sigma_2^2 + 2\theta(1 - \theta)\sigma_{12}.
\]

\begin{definition}[Mean--variance frontier]\label{def:b2:asset-prices:frontier}
The \emph{mean--variance frontier}\index{mean--variance frontier} is the set of portfolios with the
smallest variance for each mean return. Its upper branch, above the portfolio of minimum variance, is
the set an investor who likes mean and dislikes variance can choose from.
\end{definition}

Add a risk-free asset. Any mix of $R_f$ and a risky portfolio $t$ lies on the straight line from $(0,
R_f)$ through $(\sigma_t, \mathbb{E}[\tilde R_t])$ in the plane of standard deviation and mean, with
slope $(\mathbb{E}[\tilde R_t] - R_f)/\sigma_t$, the portfolio's \emph{Sharpe ratio}. The steepest such
line touches the frontier at the \emph{tangency portfolio}; every mean--variance investor holds the
safe asset and the tangency portfolio, in proportions set by its aversion to variance.

\begin{omfigure}
\begin{tikzpicture}[econ]
\begin{axis}[omrecord, grid=none, xmin=0, xmax=20, ymin=0, ymax=8.5, xlabel={Standard deviation
    (per cent)}, ylabel={Mean return (per cent)}, width=0.66\linewidth, height=0.5\linewidth]
  \addplot[black, very thick, domain=-0.25:1.15, samples=80, variable=\t]
    ({100*sqrt(\t^2*0.0324 + (1-\t)^2*0.0036 + 2*\t*(1-\t)*0.00216)}, {100*(0.07*\t + 0.03*(1-\t))});
  \addplot[omDef, thick, domain=0:18, samples=2] {2 + 0.30003*x};
  \node[eqpoint] at (axis cs:18,7) {};
  \node[anchor=north west, font=\footnotesize] at (axis cs:18,7) {equity};
  \node[eqpoint] at (axis cs:6,3) {};
  \node[anchor=north west, font=\footnotesize] at (axis cs:6,3) {bonds};
  \node[eqpoint] at (axis cs:8.973,4.692) {};
  \node[anchor=north west, font=\footnotesize] at (axis cs:8.973,4.692) {tangency $t$};
  \node[eqpoint] at (axis cs:0,2) {};
  \node[anchor=west, font=\footnotesize] at (axis cs:0.3,1.6) {$R_f$};
  \node[anchor=south east, font=\footnotesize, omDef] at (axis cs:15,6.8) {capital line, slope 0.30};
\end{axis}
\end{tikzpicture}

{\small The \omterm{def:b2:asset-prices:frontier}{mean--variance frontier} of equity (mean 7, standard deviation 18) and bonds (3 and 6),
correlation 0.2, with a safe rate of 2 per cent: the weekend problem's numbers. The tangency portfolio
holds 42.3 per cent equity; the capital line through it has the highest Sharpe ratio, $0.300$.}
\end{omfigure}

\begin{definition}[Beta of an asset]\label{def:b2:asset-prices:beta}
The \emph{beta}\index{beta (of an asset)} of asset $i$ with respect to a portfolio $t$ is
$\beta_i = \mathrm{cov}(\tilde R_i, \tilde R_t)/\sigma_t^2$: the slope of a regression of the asset's
return on the portfolio's.
\end{definition}

\begin{proposition}[Expected returns on the frontier]\label{prop:b2:asset-prices:beta-pricing}
If $t$ is the tangency portfolio, every asset satisfies
\[
\mathbb{E}[\tilde R_i] - R_f = \beta_i\,\bigl(\mathbb{E}[\tilde R_t] - R_f\bigr).
\]
\end{proposition}

\begin{proof}
The tangency portfolio maximises the Sharpe ratio over all portfolios of risky assets. Move a small
weight $\varepsilon$ from $t$ into asset $i$, financed at $R_f$: the mean \omterm{def:b2:asset-prices:excess-return}{excess return} changes by
$\varepsilon(\mathbb{E}[\tilde R_i] - R_f)$ and the variance by $2\varepsilon\,\mathrm{cov}(\tilde R_i,
\tilde R_t)$, to first order, so the standard deviation by $\varepsilon\,\mathrm{cov}(\tilde R_i, \tilde
R_t)/\sigma_t$. At a maximum of the Sharpe ratio the ratio of the two changes equals the ratio itself:
$(\mathbb{E}[\tilde R_i] - R_f)\sigma_t/\mathrm{cov}(\tilde R_i, \tilde R_t) = (\mathbb{E}[\tilde R_t]
- R_f)/\sigma_t$. Rearranging gives the formula.
\end{proof}

\begin{definition}[Capital asset pricing model]\label{def:b2:asset-prices:capm}
The \emph{capital asset pricing model}\index{capital asset pricing model} (CAPM) adds that the tangency
portfolio is the market portfolio of all risky assets, as it is when every investor chooses by mean and
variance with the same beliefs: then $\mathbb{E}[\tilde R_i] - R_f = \beta_i(\mathbb{E}[\tilde R_m] -
R_f)$, with \omterm{def:b2:asset-prices:beta}{betas} measured against the market. Plotted against \omterm{def:b2:asset-prices:beta}{beta}, expected returns lie on a line,
the \emph{security market line}.
\end{definition}

\begin{omfigure}
\begin{tikzpicture}[econ]
\begin{axis}[omrecord, grid=none, xmin=0, xmax=2, ymin=0, ymax=14, xtick={0,1,2},
    ytick={2,8}, yticklabels={$R_f$, $\mathbb{E}[\tilde R_m]$}, xlabel={Beta},
    ylabel={Mean return}, width=0.6\linewidth, height=0.45\linewidth]
  \addplot[omDef, very thick, domain=0:2, samples=2] {2 + 6*x};
  \addplot[omThm, thick, dashed, domain=0:2, samples=2] {6.5 + 1.5*x};
  \draw[guide] (axis cs:1,0) -- (axis cs:1,8) -- (axis cs:0,8);
  \node[eqpoint] at (axis cs:1,8) {};
  \node[anchor=south east, font=\footnotesize, omDef] at (axis cs:1.5,11) {CAPM};
  \node[anchor=north west, font=\footnotesize, omThm] at (axis cs:1.3,8.2) {a flatter line};
\end{axis}
\end{tikzpicture}

{\small The security market line of the CAPM (solid) and a flatter line of the kind the tests report
(dashed): \omterm{def:b2:asset-prices:beta}{low-beta} assets earn more, and \omterm{def:b2:asset-prices:beta}{high-beta} assets less, than the CAPM predicts. An illustration
with round numbers, not data.}
\end{omfigure}

\begin{model}{Capital asset pricing model}
\assumes Investors choose portfolios by the mean and variance of one period's return, agree on the
distribution of returns, and borrow and lend at one safe rate.
\predicts Expected \omterm{def:b2:asset-prices:excess-return}{excess returns} proportional to market \omterm{def:b2:asset-prices:beta}{betas}: a security market line through $R_f$
with slope equal to the market's premium; nothing else priced.
\tested The cross-section of average returns on American shares sorted by \omterm{def:b2:asset-prices:beta}{beta} and by size (Fama and
French 1992).
\breaks Once \omterm{def:b2:asset-prices:beta}{beta} is allowed to vary apart from size, the relation between \omterm{def:b2:asset-prices:beta}{beta} and average return is
% ledger: b2:ff1992-beta
flat, even with \omterm{def:b2:asset-prices:beta}{beta} as the only regressor; size and book-to-market capture the cross-section of
average returns (Fama and French 1992).
\end{model}

The chart below sets each country's three assets in the plane of standard deviation and mean, over
1950--2020. Within every country, the asset with the more volatile return paid more on average. Across
countries, a more volatile stock market did not pay more: the correlation between the standard
% figdata: summary (eq_sd_mean_corr 0.01)
deviation and the mean of the 15 equity markets is 0.01. By \cref{prop:b2:asset-prices:covariance}
that is no surprise, since what is priced is covariance with what investors fear, not volatility.

\begin{omchart}{Mean against standard deviation of annual real returns, 1950--2020, in per cent: one
  point per country and asset, 15 countries with complete series.}{Jord\`a--Schularick--Taylor
  Macrohistory Database}
\begin{tikzpicture}
\begin{axis}[omscatter, xmin=0, xmax=35, ymin=-2, ymax=14, xlabel={Standard deviation (per cent)},
    ylabel={Mean real return (per cent)}, legend pos=north west, height=0.5\linewidth,
    width=0.75\linewidth]
  \addplot+[only marks, mark=*, omDef, mark size=1.6pt]
    table[x=equity_sd, y=equity_mean] {figdata/bachelor-2/04-asset-prices/out/riskreturn_wide.csv};
  \addplot+[only marks, mark=square*, omThm, mark size=1.6pt]
    table[x=bond_sd, y=bond_mean] {figdata/bachelor-2/04-asset-prices/out/riskreturn_wide.csv};
  \addplot+[only marks, mark=triangle*, black!60, mark size=1.8pt]
    table[x=bill_sd, y=bill_mean] {figdata/bachelor-2/04-asset-prices/out/riskreturn_wide.csv};
  \legend{Equities, Bonds, Bills}
\end{axis}
\end{tikzpicture}
\end{omchart}

\section{Efficient markets}\label{sec:b2:asset-prices:efficient}

\begin{definition}[Efficient-markets hypothesis]\label{def:b2:asset-prices:emh}
The \emph{efficient-markets hypothesis}\index{efficient-markets hypothesis} says that prices reflect a
given set of information fully, so that no strategy based on that information earns more than the
return that compensates its risk. Its \emph{weak form} takes the information to be past prices; the
\emph{semi-strong form}, all public information; the \emph{strong form}, all information, private
included.
\end{definition}

\begin{proposition}[Unpredictable excess returns]\label{prop:b2:asset-prices:unpredictable}
If expected returns are constant and prices reflect past returns, the \omterm{def:b2:asset-prices:excess-return}{excess return} of a year is
uncorrelated with every past return; in particular its first-order autocorrelation is zero.
\end{proposition}

\begin{proof}
If $\mathbb{E}[\tilde R_{t+1} \mid \text{past returns}] = \bar R$, a constant, then
$\mathrm{cov}(\tilde R_{t+1}, R_{t-k}) = \mathbb{E}[(R_{t-k} - \mathbb{E}R)(\mathbb{E}[\tilde R_{t+1}
\mid \text{past}] - \bar R)] = 0$ for every $k \ge 0$, by the law of iterated expectations.
\end{proof}

The weak form is easy to test with the JST data. A sample autocorrelation from $n$ independent years
has a standard error of about $1/\sqrt n$; with 71 years, values beyond $2/\sqrt{71} = 0.237$ are
% figdata: summary (band 0.237, n_outside_band 2)
unlikely if there is no predictability. Two of the 15 countries lie beyond the band, Japan and
Portugal, against about one expected by chance at that level.

\begin{omchart}{First-order autocorrelation of annual real equity returns, 1950--2020, by country. The
  dotted lines are $\pm 2/\sqrt{71}$, the band within which about 95 per cent of sample
  autocorrelations fall when returns are unpredictable.}{Jord\`a--Schularick--Taylor Macrohistory
  Database}
\begin{tikzpicture}
\begin{axis}[omrecord, ybar, bar width=9pt, xmin=0.3, xmax=15.7, ymin=-0.35, ymax=0.35,
    xtick=data, xticklabels from table={figdata/bachelor-2/04-asset-prices/out/autocorr.csv}{iso3},
    x tick label style={font=\scriptsize, rotate=90}, ylabel={Autocorrelation}, grid=none,
    ymajorgrids, height=0.42\linewidth, width=0.85\linewidth]
  \addplot[fill=omDef!70, draw=omDef] table[x=pos, y=rho1]
    {figdata/bachelor-2/04-asset-prices/out/autocorr.csv};
  \draw[black!70, densely dotted, thick] (axis cs:0.3,0.237) -- (axis cs:15.7,0.237);
  \draw[black!70, densely dotted, thick] (axis cs:0.3,-0.237) -- (axis cs:15.7,-0.237);
\end{axis}
\end{tikzpicture}
\end{omchart}

The semi-strong form is tested by \emph{event studies}: around a public announcement (a merger, an
earnings report, a change in a dividend), the share price should jump once, at the announcement, and
drift neither before (unless information leaked) nor after. The test measures the \emph{abnormal
return}, the return minus what a pricing model predicts, day by day around the event.

\begin{definition}[Anomaly]\label{def:b2:asset-prices:anomaly}
An \emph{anomaly}\index{anomaly (asset pricing)} is a pattern in average returns that a given pricing
model does not explain: a characteristic of assets, known in advance, that predicts their returns once
their risk, as the model measures it, is accounted for.
\end{definition}

Every test of efficiency needs a model of the return that compensates risk, so it tests the two
together: the \emph{joint-hypothesis problem}. An \omterm{def:b2:asset-prices:anomaly}{anomaly} may reveal prices that ignore information, or
a model that measures risk badly. Size and book-to-market (Fama and French 1992) are anomalies for the
CAPM; added to the pricing model as risk factors, they stop being anomalies for the new model. The
sharpest test of an \omterm{def:b2:asset-prices:anomaly}{anomaly} uses data collected after it was reported: a pattern found by searching many
candidates on one sample should vanish there, and one that reflects risk or a lasting feature of
markets should not. The quant series of this collection treats the factor models and their tests at
length.

\begin{debate}[Are stock prices too volatile?]
If a share's price is the expected present value of its dividends at a constant discount rate, the
price must vary less than the realised present value of the dividends that followed, since a forecast
varies less than the thing forecast. Shiller (1981) derived such bounds and compared them with American
% ledger: b2:shiller1981, b2:mm1986, b2:cs1988
data: prices moved much more than the dividends that followed seemed to justify. Marsh and Merton
(1986) showed that the bounds can be violated by efficient prices if managers smooth dividends. Campbell
and Shiller (1988) linked the dividend--price ratio to future dividend growth and future discount rates,
so that the question became how much of the movement in prices comes from changing discount rates,
that is from changing expected returns. The evidence is clear that the dividend--price ratio moves; it
is still open how much of its movement reflects rational changes in required returns and how much
reflects prices straying from fundamentals.
\end{debate}

\begin{model}{Efficient markets}
\assumes Prices reflect all information of a given kind; expected returns are those that compensate
risk, as measured by a pricing model.
\predicts No strategy based on past prices (weak form) or public information (semi-strong form) beats
the risk-adjusted return; prices jump at news and do not drift after it.
\tested Autocorrelations of annual returns (this chapter's chart: close to zero in most countries);
event studies.
\breaks Every test is joint with a pricing model, and the pricing model of the day leaves anomalies
(size and book-to-market for the CAPM, Fama and French 1992); whether prices vary more than
fundamentals allow is still debated (above).
\end{model}

\section{Exercises}\label{sec:b2:asset-prices:exercises}

\begin{exercise}[$\star$]\label{exo:b2:asset-prices:1}
Two states. A bond pays $1$ in both and costs $0.96$; a share pays $2$ in the first state and $0.5$ in
the second and costs $1.1$. Find the \omterm{def:b2:asset-prices:state-price}{state prices}. Is there an arbitrage?
\end{exercise}

\begin{exercise}[$\star$]\label{exo:b2:asset-prices:2}
A saver has $\beta = 0.98$, CRRA utility with $\gamma = 2$, and expects consumption to grow by 2 per
cent with certainty. Find the gross risk-free return. Compare it with the mean real return on bills in
the JST data (about 1 per cent).
\end{exercise}

\begin{exercise}[$\star$]\label{exo:b2:asset-prices:3}
Read the opening chart: which asset did best in each country, and roughly by what factor did an
investment in equities grow in France and in the United States? The final values are 11.0 and 177.6;
convert each into an average annual real return over the 71 years.
\end{exercise}

\begin{exercise}[$\star$]\label{exo:b2:asset-prices:4}
Name an observation that would refute the weak form of the \omterm{def:b2:asset-prices:emh}{efficient-markets hypothesis}, and say why
a single year in which a trading rule beat the market would not.
\end{exercise}

\begin{exercise}[$\star\star$]\label{exo:b2:asset-prices:5}
A share's return has a covariance of $0.024$ with the market's, whose variance is $0.04$. The safe rate
is 1 per cent and the market's expected return 7 per cent. Find the share's \omterm{def:b2:asset-prices:beta}{beta} and its expected
return under the CAPM.
\end{exercise}

\begin{exercise}[$\star\star$]\label{exo:b2:asset-prices:6}
Consumption growth has a standard deviation of 2 per cent, an equity index 16 per cent, and their
correlation is $0.5$. Using $\mathbb{E}[\tilde R] - R_f \approx \gamma\,\mathrm{cov}(\tilde g, \tilde R)$,
predict the premium for $\gamma = 2$ and $\gamma = 10$. What $\gamma$ gives a premium of 6 per cent?
\end{exercise}

\begin{exercise}[$\star\star$]\label{exo:b2:asset-prices:7}
Data and Python. With \texttt{jst(isos=["USA"], since=1899, until=2020, cols=["eq\_tr", "cpi"])} from
\texttt{econdata}, compute the real equity return of each year 1900--2020 for the United States and the
United Kingdom, then the arithmetic mean and the geometric mean. Why is the geometric mean smaller, and
which one tells an investor what a unit invested in 1899 became?
\end{exercise}

\begin{exercise}[$\star\star$]\label{exo:b2:asset-prices:8}
Two uncorrelated assets have standard deviations of 20 and 10 per cent. Find the share of the first in
the portfolio of minimum variance, and its standard deviation.
\end{exercise}

\begin{exercise}[$\star\star$]\label{exo:b2:asset-prices:9}
A study finds that shares of small firms earned 3 points a year more than the CAPM predicts. Give two
explanations that the joint-hypothesis problem leaves open, and an observation that would favour each.
\end{exercise}

\begin{exercise}[$\star\star\star$]\label{exo:b2:asset-prices:10}
Prove the converse in \cref{prop:b2:asset-prices:state-prices}: if there is no arbitrage, both \omterm{def:b2:asset-prices:state-price}{state
prices} are positive. Then show by an example with a single security, paying $1$ in state 1 and $0$ in
state 2 at price $0.5$, that without a second security the \omterm{def:b2:asset-prices:state-price}{state price} of state 2 is not determined.
\end{exercise}

\begin{exercise}[$\star\star\star$]\label{exo:b2:asset-prices:11}
Let $\tilde\zeta = \beta\tilde g^{-\gamma}$, with $\ln\tilde g$ and $\ln\tilde R$ jointly normal. Using
$\mathbb{E}[e^{\tilde y}] = e^{\mathbb{E}\tilde y + \mathrm{var}(\tilde y)/2}$ for a normal $\tilde y$,
show that $\ln\mathbb{E}[\tilde R] - \ln R_f = \gamma\,\mathrm{cov}(\ln\tilde g, \ln\tilde R)$.
\end{exercise}

\begin{exercise}[$\star\star\star$]\label{exo:b2:asset-prices:12}
Using the listing's output: (a) explain why $\gamma \ge (\mathbb{E}[\tilde\xi]/\sigma_R)/\sigma_g$ and
check the value 25.1. (b) A household with CRRA utility and $\gamma = 25$ faces a fifty-fifty gamble of
losing or gaining 20 per cent of its wealth. Find its \omterm{def:b2:expected-utility:risk-premium}{risk premium} as a share of wealth. (c) What does
the answer say about the size of $\gamma$ the puzzle needs?
\end{exercise}

\section{Problem: The Brenmoor pension fund}\label{sec:b2:asset-prices:problem}

\begin{problem}[{Weekend problem --- the pension fund of Brenmoor's municipal workers, in Brenland, an
invented economy, prices its assets, finds the discount factor its members' consumption implies,
chooses a mix of equities and bonds, and tests whether past returns predict future ones}]
\label{pb:b2:asset-prices:1}
Next year the Brenland economy is in a boom or a slump, each with probability $0.5$. A Brenland
government bond costs 100 dalens and pays 105 in both states; a unit of the Brenland equity index costs
95 and pays 130 in a boom and 80 in a slump. Round \omterm{def:b2:asset-prices:state-price}{state prices} and \omterm{def:b2:intertemporal-choice:discount-factor}{discount factors} to four decimals,
returns to two decimals of a per cent.

\medskip
\noindent\textbf{Part I --- No arbitrage and state prices.}
\begin{enumerate}
  \item Write the payoff table. Are the two payoffs proportional?
  \item Find the \omterm{def:b2:asset-prices:state-price}{state prices} $\psi_1$ (boom) and $\psi_2$ (slump).
  \item Do these prices admit an arbitrage?
  \item Find the gross risk-free return.
  \item An option pays 30 in a boom and nothing in a slump. Find its no-arbitrage price.
  \item A dealer sells the option at 10. Build a portfolio of equity and bonds with the option's
        payoffs, and an arbitrage that exploits the dealer's price. What does it earn today?
\end{enumerate}

\medskip
\noindent\textbf{Part II --- The discount factor.}
\begin{enumerate}[resume]
  \item Find the \omterm{def:b2:asset-prices:sdf}{stochastic discount factor} in each state.
  \item Check that $\mathbb{E}[\tilde\zeta\tilde R] = 1$ for the equity.
  \item Find the equity's expected gross return and its premium over the bond.
  \item Compute $\mathrm{cov}(\tilde\zeta, \tilde R)$ and check \cref{prop:b2:asset-prices:covariance}.
  \item The fund's members consume 4 per cent more next year in a boom and 2 per cent less in a slump.
        Find the $\gamma$ and $\beta$ for which $\tilde\zeta = \beta\tilde g^{-\gamma}$.
  \item Is the $\gamma$ of question 11 plausible? Compare with \cref{exo:b2:asset-prices:12}.
\end{enumerate}

\medskip
\noindent\textbf{Part III --- Mean--variance choice.} Over many years, the fund expects equity to
return 7 per cent with a standard deviation of 18, and bonds 3 per cent with a standard deviation of
6; the correlation is $0.2$. The fund chooses its equity share $\theta$ to maximise
$\mathbb{E}[\tilde R_\theta] - \tfrac{\gamma}{2}\sigma^2_\theta$.
\begin{enumerate}[resume]
  \item Find the mean and the standard deviation of a 60/40 portfolio (60 per cent equity).
  \item Find the portfolio of minimum variance and its standard deviation.
  \item Find the optimal equity share as a function of $\gamma$, and its value for $\gamma = 2$, $3$ and
        $5$.
  \item Which $\gamma$ makes the 60/40 portfolio optimal?
  \item A safe asset now returns 2 per cent. Find the tangency portfolio, its mean, standard deviation
        and Sharpe ratio.
  \item Compare the Sharpe ratio of the 60/40 portfolio with the tangency portfolio's.
  \item With the safe asset, how much should a fund with $\gamma = 3$ put in the tangency portfolio,
        and what is its equity share overall?
\end{enumerate}

\medskip
\noindent\textbf{Part IV --- Do past returns predict?} The fund's analyst has 60 years of Brenland
equity returns (\texttt{brenland\_returns.csv}, invented).
\begin{enumerate}[resume]
  \item The mean equity return is 5.91 per cent and the standard deviation 17.55. How precise is the
        mean (its standard error)?
  \item The first-order autocorrelation is $-0.153$. Test whether it is zero at the 5 per cent level.
  \item In the 38 years after a rise, equity returned 6.28 per cent on average; in the 21 years after a
        fall, 5.12. Test whether the difference is zero.
  \item The returns were drawn independently, year by year. If the analyst had tried twenty different
        rules on these data, how many would you expect to look significant at the 5 per cent level?
  \item Why does efficiency not require returns to be unpredictable when expected returns change over
        time?
  \item State the result: the fund's optimal equity share for $\gamma = 3$ without a safe asset, the
        $\gamma$ that justifies a 60/40 portfolio, and what the predictability test finds.
\end{enumerate}
\end{problem}
